\documentclass[a4paper, 11pt]{article}
    \usepackage{vntex}
    %\usepackage[english,vietnam]{babel}
    %\usepackage[utf8]{inputenc}
    
    %\usepackage[utf8]{inputenc}
    %\usepackage[francais]{babel}
    \usepackage{a4wide,amssymb,epsfig,latexsym,multicol,array,hhline,fancyhdr}
    
    \usepackage{amsmath}
    \usepackage{lastpage}
    \usepackage[lined,boxed,commentsnumbered]{algorithm2e}
    \usepackage{enumerate}
    \usepackage{color}
    \usepackage{graphicx, listings}							% Standard graphics package
    \usepackage{array}
    \usepackage{tabularx, caption}
    \usepackage{multirow}
    \usepackage{multicol}
    \usepackage{rotating}
    \usepackage{graphics}
    \usepackage{geometry}
    \usepackage{setspace}
    \usepackage{epsfig}
    \usepackage{tikz}
    \usetikzlibrary{arrows,snakes,backgrounds}
    \usepackage{hyperref}
    \hypersetup{urlcolor=blue,linkcolor=black,citecolor=black,colorlinks=true} 
    %\usepackage{pstcol} 								% PSTricks with the standard color package
    
    \newtheorem{theorem}{{\bf Định lý}}
    \newtheorem{property}{{\bf Tính chất}}
    \newtheorem{proposition}{{\bf Mệnh đề}}
    \newtheorem{corollary}[proposition]{{\bf Hệ quả}}
    \newtheorem{lemma}[proposition]{{\bf Bổ đề}}
    
    \definecolor{codegreen}{rgb}{0,0.6,0}
    \definecolor{codegray}{rgb}{0.5,0.5,0.5}
    \definecolor{codepurple}{rgb}{0.58,0,0.82}
    \definecolor{backcolour}{rgb}{0.95,0.95,0.92}
    \lstdefinestyle{mystyle}{
        backgroundcolor=\color{backcolour},
        commentstyle=\color{codegreen},
        keywordstyle=\color{magenta},
        numberstyle=\tiny\color{codegray},
        stringstyle=\color{codepurple},
        basicstyle=\ttfamily\small,
        breakatwhitespace=false,
        breaklines=false,
        captionpos=b,
        keepspaces=true,
        numbers=left,
        numbersep=5pt,
        showspaces=false,
        showstringspaces=false,
        showtabs=false,
        tabsize=2
    }
    \lstset{style=mystyle} 
    %\usepackage{fancyhdr}
    \setlength{\headheight}{40pt}
    \pagestyle{fancy}
    \fancyhead{} % clear all header fields
    \fancyhead[L]{
     \begin{tabular}{rl}
        \begin{picture}(25,15)(0,0)
        \put(0,-8){\includegraphics[width=8mm, height=8mm]{hcmut.png}}
        %\put(0,-8){\epsfig{width=10mm,figure=hcmut.eps}}
       \end{picture}&
        %\includegraphics[width=8mm, height=8mm]{hcmut.png} & %
        \begin{tabular}{l}
            \textbf{\bf \ttfamily Trường Đại Học Bách Khoa Tp.Hồ Chí Minh}\\
            \textbf{\bf \ttfamily Khoa Khoa Học và Kỹ Thuật Máy Tính}
        \end{tabular} 	
     \end{tabular}
    }
    \fancyhead[R]{
        \begin{tabular}{l}
            \tiny \bf \\
            \tiny \bf 
        \end{tabular}  }
    \fancyfoot{} % clear all footer fields
    \fancyfoot[L]{\scriptsize \ttfamily Bài tập lớn Nhập môn trí tuệ nhân tạo 1}
    \fancyfoot[R]{\scriptsize \ttfamily Trang {\thepage}/\pageref{LastPage}}
    \renewcommand{\headrulewidth}{0.3pt}
    \renewcommand{\footrulewidth}{0.3pt}
    
    
    %%%
    \setcounter{secnumdepth}{4}
    \setcounter{tocdepth}{3}
    \makeatletter
    \newcounter {subsubsubsection}[subsubsection]
    \renewcommand\thesubsubsubsection{\thesubsubsection .\@alph\c@subsubsubsection}
    \newcommand\subsubsubsection{\@startsection{subsubsubsection}{4}{\z@}%
                                         {-3.25ex\@plus -1ex \@minus -.2ex}%
                                         {1.5ex \@plus .2ex}%
                                         {\normalfont\normalsize\bfseries}}
    \newcommand*\l@subsubsubsection{\@dottedtocline{3}{10.0em}{4.1em}}
    \newcommand*{\subsubsubsectionmark}[1]{}
    \makeatother
    
    
\begin{document}

\begin{titlepage}
	\begin{center}
		ĐẠI HỌC QUỐC GIA THÀNH PHỐ HỒ CHÍ MINH \\
		TRƯỜNG ĐẠI HỌC BÁCH KHOA \\
		KHOA KHOA HỌC - KỸ THUẬT MÁY TÍNH
	\end{center}

	\vspace{1cm}

	\begin{figure}[h!]
		\begin{center}
			\includegraphics[width=3cm]{hcmut.png}
		\end{center}
	\end{figure}

	\vspace{1cm}


	\begin{center}
		\begin{tabular}{c}
			\multicolumn{1}{l}{\textbf{{\Large NHẬP MÔN TRÍ TUỆ NHÂN TẠO}}} \\
			~~                                                              \\
			\hline
			\\
			\multicolumn{1}{l}{\textbf{{\Large Bài tập lớn 1}}}             \\
			\\
			\textbf{{\Huge TÌM KIẾM (SEARCHING)}}                           \\
			\\
			\hline
		\end{tabular}
	\end{center}

	\vspace{3cm}

	\begin{table}[h]
		\begin{tabular}{rrl}
			\hspace{5 cm} & Giảng viên:     & Vương Bá Thịnh                   \\
			              & Nhóm sinh viên: & Nguyễn Lê Huy - 1611290          \\
			              &                 & Trương Hoàng Huy - 1611352       \\
			              &                 & Nguyễn Anh Khoa - 1611617        \\
			              &                 & Nguyễn Quang Hoàng Lâm - 1611743 \\
			              &                 & Trần Đăng Khôi - 1611660         \\
		\end{tabular}
	\end{table}

	\begin{center}
		{\footnotesize TP. HỒ CHÍ MINH, THÁNG 5/2018}
	\end{center}
\end{titlepage}


%\thispagestyle{empty}

\newpage
\tableofcontents
\newpage
\section{Đề bài}
Sinh viên phải hiện thực ba giải thuật: Depth-first search, Breadth-first search,
và Simple Hill Climbing (hoặc 1 giải thuật heuristic khác) để giải 33 màn của trò chơi
\textbf{Bloxorz}.
\section{Cài đặt}
Dự án này dùng console, đã được kiểm tra trên linux, còn trên window thì không
đảm bảo hoạt động ổn định.
\subsection{Yêu cầu}
Để có thể chạy được các file python, ta cần phải có đủ các yếu tố sau:
\begin{itemize}
	\item Phải có python 3.6 trở lên
	\item Có pip để cài đặt các packages cần thiết
\end{itemize}
\subsection{Cách cài đặt và chạy}
\begin{lstlisting}
> pip install -r requirements.txt
> python main.py
\end{lstlisting}
Hoặc ta có thể dùng virtualenv
\begin{lstlisting}
> virtualenv bloxorz-env
> source bloxorz-env/bin/activate
> pip install -r requirements.txt
> python main.py    
\end{lstlisting}
Ta chạy trò chơi bằng câu lệnh "\texttt{python main.py}" hoặc "\texttt{python3 main.py}"
\newpage
\section{Cấu trúc files}
\begin{lstlisting}
    /------------------------
    |---main.py
    |---raw/
    |---|---stage1
    |---|---stage2
    |---|---...
    |---bloxorz/
    |---|---common/
    |---|---|---getKey.py
    |---|---|---menu.py
    |---|---|---moves.py
    |---|---game/
    |---|---|---play.py
    |---|---|---Stage.py
    |---|---|---Tile.py
    |---|---|---TileType.py
    |---|---solver/
    |---|---|---Solver.py
    |---|---|---State.py
    |---|---|---Block.py
    |---|---|---mode.py
    |---|---|
    |---|---|---algorithm implement files
    |---|---stages/
    |---|---|---GenStage.py
    |---|---|---stage1.py
    |---|---|---stage2.py
    |---|---...
\end{lstlisting}
\section{Cấu trúc trò chơi}
\subsection{Tiles (Các lát gạch)}
\subsubsection{Giải thích}
Mỗi lát gạch là một ô trong trò chơi. Có nhiều loại gạch là:
\begin{enumerate}
	\item Lát gạch bình thường (tile): Loại này chỉ là một ô bình thường.
	\item Lát gạch mềm (soft tile): Loại này không thể giữ 2 khối cùng 1 lúc, tức là 2 khối lập
	      phương chồng lên nhau cùng đứng trên ô gạch này thì sẽ làm vỡ gạch và rớt xuống.
	\item Lát gạch "nút" (button tile): Loại này như là nút, khi khối lập phương đè lên nó thì
	      nó sẽ kích hoạt một cái gì đó tùy thuộc vào màn chơi. Và ô gạch này cũng có nhiều loại.
	\item Lát gạch "cầu" (bridge tile): Loại này là cái sẽ bị kích hoạt bởi ô gạch "nút". Nó có thể
	      được đóng hoặc mở hoặc cả hai.
	\item Lát gạch "mục tiêu" (goal tile): Khi cả 2 khối lập phương cùng đứng trên ô gạch này là
	      sẽ chiến thắng màn chơi.
	\item Lát gạch "phân chia" (split tile): Khi cả 2 khối lập phương cùng đứng trên ô gạch này thì
	      sẽ bị phân chia thành 2 khối lập phương tách biệt nhau ở 2 vị trí khác nhau.
\end{enumerate}
Lát gạch "nút" được chia thành 2 loại chính là \textbf{soft} và \textbf{hard}, \textbf{soft} chỉ cần
1 khối lập phương đè lên còn \textbf{hard} thì cần cả 2 khối đè lên để kích hoạt.
Mỗi loại đó thì được chia thành các loại sau:
\begin{enumerate}
	\item Loại thường: Loại này thì có nhiệm vụ toggle, tức là đóng và mở.
	\item Special: Loại này chỉ có đóng hoặc mở chứ không toggle.
	\item Hell: Loại này kết hợp cả loại thường và special.
\end{enumerate}
Sau khi xem xét các loại trên, cuối cùng ta có được một lớp enum như sau:
\begin{lstlisting}
// c++
enum class TileType : int {
    normal,
    goal,
    split,
    soft_ground,

    soft_button,
    hard_button,

    soft_special_button,
    hard_special_button

    soft_hell_button,
    hard_hell_button
};
\end{lstlisting}
\subsubsection{Lớp Tile}
Lớp Tile được hiện thực như sau:
\begin{lstlisting}
// c++
class Tile {
    public:
    TileType type;
    bool valid;         // if standing on this tile is ok
    int split_place[4]; // place after split
    
    Tile*[] toggle;       // Tiles to toggle
    Tile*[] open;         // Tiles to change to open
    Tile*[] close;        // Tiles to change to close
    
    int* trigger(bool standing);
};
\end{lstlisting}
Do có nhiều loại gạch nên ta có nhiều cách construct:\\~\\
Loại gạch thường là dễ nhất: \texttt{Tile()}\\~\\
Loại gạch "mục tiêu" và "mềm" như sau: \\\texttt{Tile(TileType.goal)} và
\texttt{Tile(TileType.soft\_ground)}\\~\\
Loại gạch "cầu" cần phải biết lúc đầu nó có hay chưa nên có thêm biên bool như sau:
\texttt{Tile(TileType.bridge, true)} hoặc \texttt{Tile(TileType.bridge, false)}\\~\\
Loại gạch "nút" là khó nhất, và ta đã làm như sau:
\begin{lstlisting}
// c++

// making a button which will toggle on active
Tile(TileType.soft_button, Tile*[] toggle);
    
// making a button which will open gates only
Tile(TileType.soft_special_button, Tile*[] open, NULL);
    
// making a button which will close gates only
Tile(TileType.soft_special_button, NULL, Tile*[] close);
    
// making a button which will open and closes set of bridges
Tile(TileType.soft_special_button, Tile*[] open, Tile*[] close);
    
// making a button which will open and closes set of 
bridges but also toggle some
Tile(TileType.soft_hell_button, Tile*[] toggle, 
     Tile*[] open, Tile*[] close);
    
// the same for hard button
\end{lstlisting}
\subsubsection{Quá trình kích hoạt}
Một lát gạch có một hàm kích hoạt (trigger function), nó sẽ được thực hiện tùy theo loại gạch như sau:
\begin{itemize}
	\item Loại "nút": tùy theo loại nút, hàm có thể đóng/mở/toggle các cây cầu
	\item Loại "phân chia": hàm sẽ trả về địa chỉ mới của các khối lập phương
	\item Loại "mềm": hàm sẽ ném (throw) tính hiệu lỗi "Fall"
	\item Loại "cầu": nếu không hiệu lực thì hàm sẽ throw "Hidden bridge"
\end{itemize}
Ngoài các trường hợp trên thì hàm sẽ không làm gì cả.
\subsection{Nền (stage)}
Lớp nền lưu màn chơi và vị trí của các khối lập phương thành các file binary.
\begin{lstlisting}
// c++
class Stage {
    public:
    Tile[][] board;  // a 2 dimension matrix Tile
    int _x;
    int _y;
    char* name;
    
    Stage(char* name, Tile[][] b, int x, int y);
    void save(char* filename);
}
\end{lstlisting}
Lớp này được truyền vào lớp trạng thái khi chơi.
\subsection{Khối lập phương (block)}
Khối lập phương là nhân vật chính trong trò chơi, mọi thứ đều diễn ra xung quanh khối lập phương.
Để phân biệt, khi 2 khối lập phương dính vào nhau thì ta gọi là Blox, còn riêng biệt thì ta gọi là Block.
Block chính nó chỉ có 1 index, index còn lại được tính toán thông qua trạng thái (đứng, nằm dọc, nằm ngang). \\~\\
Khi hàm phân chia kích hoạt, Block sẽ chuyển trạng thái của nó thành neutral. Vì thế ta cần thêm 1 block để
đại diện cho block được phân chia. Blox được hiện thực là 1 array ở trong lớp trạng thái. Ta có một bit chọn
block nào sẽ được di chuyển khi hàm di chuyển được gọi.
\subsection{Các nước di chuyển (moves)}
Đây là lớp enum như sau:
\begin{lstlisting}
// c++
enum moves {
    nomoves,
    up,
    down,
    right,
    left,
    split,
    join,
    swap
};
\end{lstlisting}
\subsection{Trạng thái (state)}
Lớp trạng thái hơi đơn giản, nó chỉ gồm cái nền và khối lập phương và một số phương thức.
\begin{lstlisting}
    // c++
    class State {
        Tile[][] board;
        Block blox[2];
        int selection = 1;
        moves[] moves;
    
        void toggleActive() {
            if (this->selection == 2)
                this->selection = 1;
            else
                this->selection = 2
        }
    };
    
    void move(State* s, moves m) {
        void* ret = NULL;
        Block* block = &(s->blox[x->selection - 1])
        try {
            block->move(m);
    
            if (block is out of bound)
                throw outofbound
    
            // trigger the tiles
            ret = s->board[block->index()].trigger();
    
            if !(block->standing || block->spliting)
                // trigger the other block
        }
        catch(everything) {
            // reverse the move
            block->move(m.reverse())
            throw error;
        }
    
        s->moves.append(m)
    
        if (ret != NULL) {
            // ret is now a index after split
            s->blox[0]->split(ret);
            // post process
            // blox[1] = new Block(ret);
    
            s->moves.append(moves.split)
        }
    
        // try to join
        // if not split, return false
        // if not able to join, return false
        // else join and return true and add to list of moves
        if (s->join())
            s->moves.append(moves.join)
    }
\end{lstlisting}
Khi viết các giải thuật, chỉ cần gọi \texttt{move(s, direction)} là xong. Nếu khi di chuyển mà có lỗi
(ra ngoài nền, đứng trên lát gạch "mềm"), thì cứ việc bỏ qua nó.
\subsection{Lời giải (solver)}
Framework này dựa trên mẫu chiến lược, tùy theo chế độ nào mà ta chọn, nó sẽ chạy giải thuật tương ứng.
Lớp chính của solver lưu trạng thái khởi đầu, trạng thái kết thúc nếu có, và chế độ giải thuật. Hơn nữa,
khi tạo các trạng thái, trạng thái đã tồn tại có thể được tạo đi tạo lại. Để ngăn chặn điều này, lớp solver
sẽ lưu một danh sách các "kí hiệu trạng thái".\\~\\
"Kí hiệu trạng thái" là một chuỗi được định nghĩa rằng trạng thái được sinh ra là duy nhất.\\~\\
Một lời gọi tới solver sẽ trông giống như thế này:
\begin{lstlisting}
// c++
BinaryLoader* binary = new BinaryLoader();  // a fake binary loader
Stage* stage = binary->load("Filename");
State* init = new State(stage);
m = mode.bfs;  // bfs algo
Solver* problem = new Solver(init, mode);
problem->solve();
problem->printSolution();
\end{lstlisting}
\section{Các giải thuật}
Sự tiêu tốn bộ nhớ được kiểm tra bằng memory\_profiler ở hàm solve trong file Solver.py
\subsection{BFS}
\subsubsection{Giải thích}
Giải thuật được hiện thực theo mã giả sau:
\begin{verbatim}
Tạo queue với 1 phần tử là trạng thái khởi đầu
Khởi tạo dãy dấu vết
Khởi tạo bộ đếm bằng 0
Lặp cho tới khi queue rỗng:
    Nếu bộ đếm > 50000 thì trả về False
    Tăng bộ đếm lên thêm 1
    Lấy cur_state = dequeue(queue)
    Nếu cur_state là đích thì trả về True
    Lặp từng khối:
        Lặp từng nước đi:
            thử:
                sinh ra trạng thái mới là new_state
                lấy dấu vết của new_state
                nếu dấu vết không có trong dãy dấu vết thì:
                    enqueue(new_state)
                    thêm dấu vết vào dãy dấu vết
            thấy lỗi thì cứ tiếp tục
Trả về False
\end{verbatim}
\subsubsection{Bảng số liệu}
Bảng số liệu về thời gian và bộ nhớ:
\begin{center}
	\begin{tabular}{|c|l|c|}
		\hline
		Màn chơi & Thời gian (s)     & Bộ nhớ (MiB) \\ \hline
		1        & 0.266571521759033 & 17.410       \\ \hline
		2        & 1.3626389503479   & 17.680       \\ \hline
		3        & 0.355221509933472 & 17.168       \\ \hline
		4        & 0.422888994216919 & 17.320       \\ \hline
		5        & 1.83321666717529  & 17.316       \\ \hline
		6        & 0.495889902114868 & 17.332       \\ \hline
		7        & 0.80048131942749  & 17.164       \\ \hline
		8        & 4.72867488861084  & 19.746       \\ \hline
		9        & 5.53647089004517  & 18.215       \\ \hline
		10       & 106.905389547348  & 22.871       \\ \hline
		11       & 1.13060832023621  & 17.219       \\ \hline
		12       & 2.10779356956482  & 17.316       \\ \hline
		13       & 1.10481071472168  & 17.484       \\ \hline
		14       & 3.52388048171997  & 17.699       \\ \hline
		15       & 81.3969686031342  & 24.523       \\ \hline
		16       & 2.33484292030334  & 17.312       \\ \hline
		17       & 5.38768458366394  & 18.148       \\ \hline
		18       & 4.20741415023804  & 17.844       \\ \hline
		19       & 1.68268489837646  & 17.121       \\ \hline
		20       & 73.6282043457031  & 25.082       \\ \hline
		21       & 1.71589279174805  & 17.441       \\ \hline
		22       & 2.1993715763092   & 17.305       \\ \hline
		23       & 78.8143999576569  & 22.977       \\ \hline
		24       & 3.81749391555786  & 17.492       \\ \hline
		25       & 3.30008316040039  & 17.328       \\ \hline
		26       & 159.813554048538  & 23.262       \\ \hline
		27       & 2.79411315917969  & 17.312       \\ \hline
		28       & 131.718682527542  & 27.406       \\ \hline
		29       & 13.9431304931641  & 18.266       \\ \hline
		30       & 5.53034996986389  & 17.633       \\ \hline
		31       & 8.13841199874878  & 17.926       \\ \hline
		32       & 4.42125082015991  & 17.477       \\ \hline
		33       & 7.48211884498596  & 19.141       \\ \hline
	\end{tabular}
\end{center}
\subsection{DFS}
\subsubsection{Giải thích}
Giải thuật được hiện thực theo mã giả sau:
\begin{verbatim}
Tạo stack là list với 1 phần tử là trạng thái khởi đầu
Khởi tạo dãy dấu vết
Khởi tạo bộ đếm bằng 0
Lặp cho tới khi stack rỗng:
    Nếu bộ đếm > 50000 thì trả về False
    Tăng bộ đếm lên thêm 1
    Lấy cur_state = phần tử cuối của stack
    Xóa phần tử cuối của stack đi
    Nếu cur_state là đích thì trả về True
    Lặp từng khối:
        Lặp từng nước đi:
            thử:
                sinh ra trạng thái mới là new_state
                lấy dấu vết của new_state
                nếu dấu vết không có trong dãy dấu vết thì:
                    thêm new_state vào list stack
                    thêm dấu vết vào dãy dấu vết
            thấy lỗi thì cứ tiếp tục
Trả về False
\end{verbatim}
\subsubsection{Bảng số liệu}
Bảng số liệu về thời gian và bộ nhớ:
\begin{center}
	\begin{tabular}{|c|l|c|}
		\hline
		Màn chơi & Thời gian (s)     & Bộ nhớ (MiB) \\ \hline
		1        & 0.135706424713135 & 17.336       \\ \hline
		2        & 1.58651113510132  & 18.691       \\ \hline
		3        & 0.313766956329346 & 17.371       \\ \hline
		4        & 0.474039077758789 & 17.320       \\ \hline
		5        & 0.696768522262573 & 18.344       \\ \hline
		6        & 0.413698673248291 & 17.285       \\ \hline
		7        & 0.701162338256836 & 17.477       \\ \hline
		8        & 0.7951500415802   & 20.508       \\ \hline
		9        & 8.05996417999268  & 27.438       \\ \hline
		10       & 11.2849471569061  & 40.230       \\ \hline
		11       & 1.04428720474243  & 17.402       \\ \hline
		12       & 2.22384285926819  & 17.883       \\ \hline
		13       & 1.14733099937439  & 18.090       \\ \hline
		14       & 1.85471248626709  & 17.871       \\ \hline
		15       & 6.61174869537354  & 24.254       \\ \hline
		16       & 4.19367742538452  & 19.793       \\ \hline
		17       & 3.84662222862244  & 19.805       \\ \hline
		18       & 2.80050277709961  & 18.938       \\ \hline
		19       & 1.08713054656982  & 17.617       \\ \hline
		20       & 64.9240102767944  & 48.371       \\ \hline
		21       & 0.979677677154541 & 17.281       \\ \hline
		22       & 1.60459566116333  & 17.855       \\ \hline
		23       & 29.9765365123749  & 51.020       \\ \hline
		24       & 0.794204235076904 & 17.324       \\ \hline
		25       & 1.6055965423584   & 17.871       \\ \hline
		26       & 69.1387176513672  & 60.266       \\ \hline
		27       & 2.36821866035461  & 17.734       \\ \hline
		28       & 120.781942844391  & 78.012       \\ \hline
		29       & 12.1229801177979  & 18.688       \\ \hline
		30       & 2.8791708946228   & 18.031       \\ \hline
		31       & 5.11534285545349  & 19.148       \\ \hline
		32       & 2.44287276268005  & 17.867       \\ \hline
		33       & 7.44981098175049  & 21.773       \\ \hline
	\end{tabular}
\end{center}
\subsection{HILL CLIMBING}
\subsubsection{Giải thích}
Giải thuật này chính xác hơn là Steepest ascent climbing. Ta giải quyết vấn đề tối ưu cục bộ bằng
backtracking. Giải thuật được hiện thực theo mã giả sau:
\begin{verbatim}
Khởi tạo dãy dấu vết
Khởi tạo bộ đếm bằng 0
Khởi tạo state là trạng thái khởi đầu
Duyệt tìm vị trí goal bằng cách dùng state
Khởi tạo stack là mảng trạng thái 1 phần tử là trạng thái ban đầu
Lặp cho đến khi stack rỗng:
    Nếu bộ đếm > 50000 thì trả về False
    Tăng bộ đếm lên thêm 1
    Gán cur_state là phần tử cuối cùng của stack
    Xóa phần tử cuối của stack đi
    Nếu cur_state là đích thì trả về True
    Khởi tạo priority là mảng rỗng
    Lặp từng khối:
        Lặp từng nước đi:
            thử:
                sinh ra trạng thái mới là new_state
                lấy dấu vết của new_state
                gán giá trị trả về của hàm EVAL(new_state, goal) là newEval 
                nếu dấu vết không có trong dãy dấu vết:
                        thêm dấu vết vào dãy dấu vết
                        thêm cặp (newEval, new_state) vào priority
            thấy lỗi thì cứ tiếp tục
    Sắp xếp priority giảm dần
    Lấy hết trạng thái trong priority bỏ vào stack theo thứ tự đó
Trả về False
\end{verbatim}
Hàm EVAL được thực hiện như sau:
\begin{verbatim}
EVAL(trạng thái, vị trí goal):
    Khởi tạo ans = 0
    Lặp từng khối lập phương:
        ans += khoảng cách giữa từng khối đến goal
    Trả về ans
\end{verbatim}
\subsubsection{Bảng số liệu}
Bảng số liệu về thời gian và bộ nhớ:
\begin{center}
	\begin{tabular}{|c|l|c|}
		\hline
		Màn chơi & Thời gian (s)       & Bộ nhớ (MiB) \\ \hline
		1        & 0.0395197868347168  & 17.305       \\ \hline
		2        & 1.7283692359924316  & 18.398       \\ \hline
		3        & 0.24228239059448242 & 17.551       \\ \hline
		4        & 0.6023290157318115  & 17.328       \\ \hline
		5        & 0.8682773113250732  & 18.129       \\ \hline
		6        & 0.5019304752349854  & 17.336       \\ \hline
		7        & 0.7346084117889404  & 17.332       \\ \hline
		8        & 0.2118821144104004  & 18.234       \\ \hline
		9        & 0.9325628280639648  & 19.246       \\ \hline
		10       & 8.855451345443726   & 31.543       \\ \hline
		11       & 1.0376548767089844  & 17.246       \\ \hline
		12       & 1.1754610538482666  & 18.195       \\ \hline
		13       & 0.46993041038513184 & 17.355       \\ \hline
		14       & 1.731593370437622   & 17.879       \\ \hline
		15       & 6.74385929107666    & 22.254       \\ \hline
		16       & 0.36205363273620605 & 17.035       \\ \hline
		17       & 4.150902032852173   & 20.875       \\ \hline
		18       & 2.7860164642333984  & 18.922       \\ \hline
		19       & 0.7718052864074707  & 17.207       \\ \hline
		20       & 74.86791396141052   & 52.602       \\ \hline
		21       & 1.0988316535949707  & 17.285       \\ \hline
		22       & 0.9383811950683594  & 18.355       \\ \hline
		23       & 29.081502676010132  & 43.809       \\ \hline
		24       & 1.2967350482940674  & 18.074       \\ \hline
		25       & 1.707021713256836   & 17.516       \\ \hline
		26       & 58.00024366378784   & 39.652       \\ \hline
		27       & 2.257538318634033   & 17.945       \\ \hline
		28       & 46.117361068725586  & 31.828       \\ \hline
		29       & 11.1941556930542    & 18.434       \\ \hline
		30       & 2.666714906692505   & 18.082       \\ \hline
		31       & 5.266216039657593   & 19.316       \\ \hline
		32       & 3.187586545944214   & 18.809       \\ \hline
		33       & 6.046768665313721   & 21.145       \\ \hline
	\end{tabular}
\end{center}
\end{document}